\clearpage\thispagestyle{empty}
\pdfbookmark[0]{Quellen zum Beitrag „Chronometrische Längenbestimmungen“}{chronometric1826-sources}
\begin{center}\textbf{Quellen zum Beitrag „Chronometrische Längenbestimmungen“}\end{center}
\noindent Diese KI-Übertragung ins heutige Deutsch folgt \textit{Chronometrische Längenbestimmungen} in \href{https://gdz.sub.uni-goettingen.de/id/PPN236007467}{\textit{Werke}, Band~VI} (Göttingen, 1874), S.~455--459. Der Erstdruck in den \textit{Astronomischen Nachrichten}, Band~V, Nr.~110, bezeichnet den Beitrag als Auszug aus einem Schreiben von Gauß an den Herausgeber. Im Exemplar der Bayerischen Staatsbibliothek steht er in den Spalten \href{https://api.digitale-sammlungen.de/iiif/image/v2/bsb10049430_00124/full/1500,/0/default.jpg}{227--228}, \href{https://api.digitale-sammlungen.de/iiif/image/v2/bsb10049430_00125/full/1500,/0/default.jpg}{229--230}, \href{https://api.digitale-sammlungen.de/iiif/image/v2/bsb10049430_00126/full/1500,/0/default.jpg}{231--232} und \href{https://api.digitale-sammlungen.de/iiif/image/v2/bsb10049430_00127/full/1500,/0/default.jpg}{233--234}. Die Werkausgabe datiert den Beitrag auf November 1826. Der gebundene Zeitschriftenband trägt die Jahreszahl 1827.
\par\smallskip\noindent Die Werkausgabe enthält Ernst Scherings Anmerkung zu den Beobachtungen von Dr. Tiarks auf der Expedition der englischen Admiralität von 1824. In ihrem Text steht die Nummer 4 zweimal, über dem $S$ unter dem ersten Wurzelzeichen steht ein Strich, und der erste Vergleich in Greenwich wird auf Juni datiert. Die Zeitschrift druckt die Nummer 3 zweimal und setzt unter das erste Wurzelzeichen ein $S$ ohne Strich. Im Fließtext nennt sie Januar, in der Tabelle Juni. Bei der Bremer Rechnung nennt sie zudem den Stand des gedachten Chronometers und am Schluss Gauß als Unterzeichner. Die Übertragung folgt dem Text und der Tabellenanordnung der Werkausgabe. KI-Übertragung und KI-Satz.
\clearpage
